% Einheit 1 von 4 – Prozentsatz berechnen (Katalog Einheit 2), Kl. 7, schwach
\clearpage
\hypertarget{e1}{}
\einheitenkopf*{Einheit 1 von 4 · Prozentsatz berechnen}
\zweigzeile{Hier lernst du auszurechnen, wie viel Prozent ein Teil vom Ganzen ist · neu in diesem Jahr · P10 oft · baut auf: Prozente als Anteile, Brüche erweitern, Runden (Kennst du schon)}
\setcounter{aufgabe}{10}

\begin{aufgabe}{Ich erkenne, was das Ganze ist. Das Ganze ist immer der ganze Streifen, also 100\,\%. Kreuze an, was hier das Ganze ist.}
\swb{\small Das Ganze $= 100\,\%$}{\streifen{100}{$0\,\%$}{$100\,\%$}}
\begin{teile}
\teil Von 30 Kindern der Klasse fahren 12 mit dem Rad. \\ \kreuz{30 Kinder}\quad\kreuz{12 Kinder}
\teil Tom hat 8 von 10 Aufgaben richtig. \\ \kreuz{8 Aufgaben}\quad\kreuz{10 Aufgaben}
\teil Ein Pullover kostete 40\,€. Jetzt ist er 10\,€ billiger. \\ \kreuz{40\,€}\quad\kreuz{10\,€}
\teil Eine Flasche fasst 1 Liter. 250\,ml sind schon getrunken. \\ \kreuz{250\,ml}\quad\kreuz{1 Liter}
\teil Im Glas sind 5 rote und 15 blaue Murmeln. Wie viel Prozent der Murmeln sind rot? \\ \kreuz{die roten Murmeln}\quad\kreuz{die blauen Murmeln}\quad\kreuz{alle Murmeln}
\end{teile}
\end{aufgabe}

\begin{aufgabe}{Ich kann einen Streifen in gleich große Teile einteilen. Teile den Streifen mit Strichen in gleich große Schritte ein.}
\swa{in 50-\%-Schritte}{\streifen[0]{0}{$0\,\%$}{$100\,\%$}}
\swa{in 25-\%-Schritte}{\streifen[0]{0}{$0\,\%$}{$100\,\%$}}
\swa{in 10-\%-Schritte}{\streifen[0]{0}{$0\,\%$}{$100\,\%$}}
\swa{in 20-\%-Schritte}{\streifen[0]{0}{$0\,\%$}{$100\,\%$}}
\end{aufgabe}

\begin{aufgabe}{Ich kann am Streifen ablesen, wie viel Prozent ein Teil ist. Ein Kästchen ist ein gleich großer Teil des Ganzen. Lies ab, wie viel Prozent gefärbt sind.}
\swb{\small Der Streifen sind 10 Kinder, gefärbt sind 4 Kinder.\\[2pt]\beispiel{1 \text{ Kästchen} &= 1 \text{ Kind} = 10\,\% \\ 4 \text{ Kästchen} &= 4 \text{ Kinder} = 40\,\%}}{\streifen{40}{$0$}{10 Kinder}}
\swz{Der Streifen sind 10 Kinder, gefärbt sind 7 Kinder.}{\streifen{70}{$0$}{10 Kinder}}
\swz{Der Streifen sind 20\,€, gefärbt sind 6\,€.}{\streifen{30}{$0\,€$}{$20\,€$}}
\swz{Der Streifen sind 4 Stunden, gefärbt ist 1 Stunde.}{\streifen[4]{25}{$0$\,h}{$4$\,h}}
\swz{Der Streifen sind 50\,m, gefärbt sind 45\,m.}{\streifen{90}{$0$\,m}{$50$\,m}}
\end{aufgabe}

\begin{aufgabe}{Ich kann ausrechnen, wie viel Prozent das sind. Schreibe den Anteil als Bruch mit dem Nenner 100 und dann in Prozent. Die Prozentzahl, die herauskommt, heißt Prozentsatz.}
\swb{\beispiel{20 \text{ von } 100 &= \tfrac{20}{100} \\ &= 20\,\%}}{\streifenfrage{20}{$0$}{$100$}}
\swz{10 von 100}{\streifenfrage{10}{$0$}{$100$}}
\swz{30 von 100}{\streifenfrage{30}{$0$}{$100$}}
\swz{50 von 100}{\streifenfrage{50}{$0$}{$100$}}
\swz{70 von 100}{\streifenfrage{70}{$0$}{$100$}}
\swfrage{Was bleibt gleich, was ändert sich?}
\end{aufgabe}

\begin{aufgabe}{Ich kann ausrechnen, wie viel Prozent das sind – weiter. Erweitere zuerst auf den Nenner 100.}
\swz{13 von 25}{\streifenfrage[0]{52}{$0$}{$25$}}
\anw{Rechne Teil : Ganzes mit dem Taschenrechner. Schreibe das Ergebnis als Dezimalzahl und dann in Prozent.}
\swz{27 von 60}{\streifenfrage[20]{45}{$0$}{$60$}}
\anw{Rechne wie eben. Runde den Prozentsatz auf eine Stelle nach dem Komma.}
\swz{39 von 47}{\streifenfrage[0]{83}{$0$}{$47$}}
\end{aufgabe}

\begin{aufgabe}{Ich kann zu einer Prozentzahl passende Zahlen finden. Finde einen Teil und ein Ganzes, die diesen Prozentsatz ergeben. Schreibe: \dots{} von \dots}
\swz[1]{50\,\%}{\streifen{50}{$0$}{$?$}}
\swz[1]{25\,\%}{\streifen[20]{25}{$0$}{$?$}}
\swz[1]{10\,\%}{\streifen{10}{$0$}{$?$}}
\swz[1]{60\,\%}{\streifen{60}{$0$}{$?$}}
\end{aufgabe}

\begin{aufgabe}{Ich kann in einem Text ausrechnen, wie viel Prozent das sind. Finde zuerst das Ganze. Rechne dann Teil : Ganzes.}
\swz{Im Schulchor singen 18 Mädchen und 12 Jungen. Wie viel Prozent der Kinder im Chor sind Mädchen?}{\streifenfrage{60}{$0$}{alle Kinder}}
\swz{Ein Rucksack kostete 50\,€. Jetzt kostet er 40\,€. Wie viel Prozent Rabatt gibt es?}{\streifenfrage{20}{$0\,€$}{$50\,€$}}
\begin{teile}
\teil Beim Klassenfest gibt es 5 Pizzen mit Pilzen und 19 Pizzen ohne Pilze. Wie viel Prozent aller Pizzen sind mit Pilzen? Runde auf eine Stelle nach dem Komma. (P10 2018)
\end{teile}
\end{aufgabe}

\begin{aufgabe}{Ich finde den Fehler in einer Prozentrechnung. Von 20 Kindern der Klasse haben 9 ein Haustier. Tim rechnet so:}
\rechnung{20 : 9 &\approx 2{,}22 = 222\,\%}
\swz{Finde den Fehler, benenne ihn und korrigiere die Rechnung.}{\streifenfrage{45}{$0$}{20 Kinder}}
\swz{Von 25 Kindern spielen 7 ein Instrument. Wie viel Prozent sind das?}{\streifenfrage[0]{28}{$0$}{25 Kinder}}
\end{aufgabe}

\begin{aufgabe}{Ich finde den Fehler in einer Prozentrechnung – weiter. In einer Tüte sind 4 Kirschbonbons und 16 andere Bonbons. Mia rechnet so:}
\rechnung{4 : 16 &= 0{,}25 = 25\,\%}
\swz{Finde den Fehler, benenne ihn und korrigiere die Rechnung.}{\streifenfrage{20}{$0$}{alle Bonbons}}
\swz{Auf einem Parkplatz stehen 9 rote und 21 andere Autos. Wie viel Prozent der Autos sind rot?}{\streifenfrage{30}{$0$}{alle Autos}}
\end{aufgabe}

\begin{aufgabe}{Ich kann begründen, warum ich durch das Ganze teile.}
\swz{Ben soll ausrechnen, wie viel Prozent 6 von 30 sind. Er fragt: „Rechne ich 6 : 30 oder 30 : 6?“ Entscheide und begründe.}{\streifen{20}{$0$}{$30$}}
\swz{Lea sagt: „Wenn ich einen Teil durch sein Ganzes teile, kommen nie mehr als 100\,\% heraus.“ Stimmt das? Begründe.}{\streifen{100}{$0\,\%$}{$100\,\%$}}
\end{aufgabe}

\begin{aufgabe}{Ich kann mit Prozenten vergleichen, wer besser trifft. Beim Basketball wirft Lea 24-mal und trifft 18-mal. Ben wirft 30-mal und trifft 21-mal. Ben sagt: „Ich habe öfter getroffen, also bin ich besser.“}
\swz{Wie viel Prozent ihrer Würfe trifft Lea?}{\streifenfrage[20]{75}{$0$}{24 Würfe}}
\swz{Wie viel Prozent seiner Würfe trifft Ben?}{\streifenfrage{70}{$0$}{30 Würfe}}
\swz{Wer trifft besser? Hat Ben recht? Begründe mit deinen Ergebnissen.}{}
\end{aufgabe}

\uebersichtskasten{Prozentsatz: Teil geteilt durch Ganzes – das ist der Anteil, dann in Prozent schreiben. \\ 14 von 40: \ $14 : 40 = 0{,}35 = 35\,\%$ \qquad 3 von 8: \ $3 : 8 = 0{,}375 = 37{,}5\,\%$}
